\setstretch{1}
\setlength{\parskip}{\baselineskip}

Usually, the first chapter of the thesis provides an introduction to the research work. Each
chapter may start with an introductory paragraph right after its title to provide some
information about its content.

\section{Section Example}
You can write in italic as follows \textit{Italic text}
You can write in bold as follows \textbf{bold text}
You can start a new paragraph by leaving an empty line as follows.

New Paragraph starts here.
\section{Another Section Example}\
There are three levels of headings in this template. Using the heading styles allows for
automatic numbering of all sections and helps in automatically generating the table of
contents.
\subsection{Sub-Section 1 Example}
\subsubsection{Sub-Sub-Section 1 Example}
You can add a new Nomenclature entry \newnom{Description}{Symbol} or \newnom{speed of light}{$c$} as follows \verb|\newnom{Description}{Symbol}|. The Nomenclature entries will be listed in the Nomenclature section in the front matter. You can also insert an entry in the Nomenclature section without appearing in the text as follows \verb|\nomenclature{Symbol}{Description}|.
